% ============================================================================
%  Método Ideal (amostragem por trem de impulsos)
%
%  Ordem dos slides:
%    A. Descrição analítica geral  (conceito, interpretação gráfica,
%       demonstração no tempo e na frequência)
%    B. Aplicação aos parâmetros fixos  (x_f(t) = cos(2 pi 2 t), f_s = 5 Hz)
%    C. Reconstrução  (filtro, frequência, tempo e verificação numérica)
%    D. Visão geral do processo + gancho para os métodos com pulso
%
%  Macros \ft... e \id... definidas no main.tex.
%  As figuras usam os próprios parâmetros fixos: f_0 = 2 Hz e f_s = 5 Hz.
% ============================================================================

% ----------------------------------------------------------------------------
%  A. DESCRIÇÃO ANALÍTICA GERAL
% ----------------------------------------------------------------------------
\begin{frame}{Método Ideal: Amostragem}{Análise no Domínio do Tempo}

    Multiplicamos o sinal $x_f(t)$ por um trem de impulsos deslocados de $T_s$.
    Cada impulso ``recolhe'' o valor instantâneo do sinal, sem qualquer retenção:
    a amostra existe apenas no instante $t = nT_s$.

    \vspace{1em}
    A equação para a modulação é dada por:
    \begin{equation*}
        x_\delta(t) = x_f(t) \cdot p(t),
        \qquad
        p(t) = \sum_{n=-\infty}^{\infty} \delta(t - nT_s)
    \end{equation*}

    Pela propriedade de peneiramento do impulso,
    $x_f(t)\,\delta(t - nT_s) = x_f(nT_s)\,\delta(t - nT_s)$:
    \begin{equation*}
        x_\delta(t) = \sum_{n=-\infty}^{\infty} x_f(nT_s)\, \delta(t - nT_s)
    \end{equation*}

    Não há pulso de retenção: o trem de impulsos é o único elemento do processo.

\end{frame}

\begin{frame}{Método Ideal: Amostragem}{Interpretação Gráfica}
\centering
\begin{minipage}[c][0.66\textheight][c]{0.92\textwidth}
\centering

% ------------------------------------------------------- (1) sinal original
\only<1>{%
\begin{tikzpicture}
    \begin{axis}[pamaxis, ylabel={$x_f(t)$}]
        \addplot[sinalcor, line width=1pt, smooth] {\sinal};
    \end{axis}
\end{tikzpicture}\par
Sinal analógico de banda base $x_f(t)$.}

% ------------------------------------------------------ (2) trem de impulsos
\only<2>{%
\begin{tikzpicture}
    \begin{axis}[pamaxis, ylabel={$p(t)$}]
        \addplot[draw=none] {0};
        \pgfplotsextra{%
            \foreach \n in {0,...,\Nam}{%
                \pgfmathsetmacro\xn{\n*\Ts}%
                \draw[sinalcor, line width=1pt, -{Stealth[length=5pt,width=4pt]}]
                    (axis cs:\xn,0) -- (axis cs:\xn,0.85);
                \draw[gray, line width=.4pt] (axis cs:\xn,-0.07) -- (axis cs:\xn,0.07);}}
        \marcaTs
    \end{axis}
\end{tikzpicture}\par
Trem de impulsos com período $T_s$, todos de área unitária.}

% ------------------------------------------- (3) produto (amostragem ideal)
\only<3>{%
\begin{tikzpicture}
    \begin{axis}[pamaxis, ylabel={$x_\delta(t)$}]
        \addplot[gray!50, dashed, line width=.7pt, smooth] {\sinal};
        \pgfplotsextra{%
            \foreach \n in {0,...,\Nam}{%
                \pgfmathsetmacro\xn{\n*\Ts}%
                \pgfmathsetmacro\yn{sin(deg(2*pi*0.75*\n*\Ts + 0.4))}%
                \draw[sinalcor, line width=1pt, -{Stealth[length=5pt,width=4pt]}]
                    (axis cs:\xn,0) -- (axis cs:\xn,\yn);}}
        \marcaTs
    \end{axis}
\end{tikzpicture}\par
Produto: impulsos de área $x_f(nT_s)$, nulos fora dos instantes de amostragem.}

\end{minipage}
\end{frame}

\begin{frame}{Método Ideal: Amostragem}{Demonstração no Domínio do Tempo}
\small
No domínio do tempo, o processo é apenas a \idamos{multiplicação do sinal por um trem
de impulsos}, sem etapa de retenção. Aplicando o peneiramento, cada impulso fica
ponderado pela sua amostra:
\begin{equation*}
  x_\delta(t) = x_f(t)\cdot \idamos{p(t)}
  = \boxed{\;\sum_{n=-\infty}^{\infty} x_f(nT_s)\, \delta(t - nT_s)\;},
  \qquad
  \idamos{p(t) = \sum_{n=-\infty}^{\infty} \delta(t - nT_s)}
\end{equation*}
\begin{center}
\begin{tikzpicture}
\begin{groupplot}[group style={group name=g, group size=3 by 1, horizontal sep=1.3cm}, ftdemot, height=0.18\textheight]
\nextgroupplot[title={(a) Sinal original\\[1pt]{\tiny\ftstrut $x_f(t)$}}]
  \addplot[ftazul!40!black, thick, domain=0:\fttmax, samples=300] {cos(deg(2*pi*\ftfzero*x))};
\nextgroupplot[title={(b) Trem de impulsos\\[1pt]{\tiny\ftstrut $p(t)$}}]
  \addplot[ycomb, ftazul, thick] coordinates {\ftimpt};
  \addplot[ftimppos, ftazul] coordinates {\ftimpt};
  \draw[gray, <->, >=stealth] (axis cs:\ftTs,1.18) -- node[above, font=\tiny, inner sep=1pt]{$T_s$} (axis cs:2*\ftTs,1.18);
\nextgroupplot[title={(c) Sinal amostrado\\[1pt]{\tiny\ftstrut $x_\delta(t)$}}]
  \addplot[gray!55, dashed, domain=0:\fttmax, samples=300] {cos(deg(2*pi*\ftfzero*x))};
  \addplot[ycomb, ftazul, thick] coordinates {\ftamostras};
  \addplot[ftimppos, ftazul] coordinates {\ftamospos};
  \addplot[ftimpneg, ftazul] coordinates {\ftamosneg};
\end{groupplot}
\ftop{1}{2}{\times}
\ftop{2}{3}{=}
\end{tikzpicture}
\end{center}
O sinal amostrado só existe nos instantes $t = nT_s$: entre as amostras não há informação.
\end{frame}

\begin{frame}{Método Ideal: Amostragem}{Demonstração no Domínio da Frequência}
\small
Aplicando a transformada de Fourier, o produto no tempo vira convolução na frequência.
Como a convolução com um impulso desloca o espectro, cada impulso de
$\idamos{P(f)}$ cria uma cópia inteira de $X_f(f)$:
\begin{equation*}
  X_\delta(f) = X_f(f) * \idamos{P(f)},
  \quad
  \idamos{P(f) = \frac{1}{T_s}\sum_{n=-\infty}^{\infty} \delta(f - nf_s)}
  \;\Longrightarrow\;
  \boxed{\; X_\delta(f) = f_s\sum_{n=-\infty}^{\infty} X_f(f - nf_s) \;}
\end{equation*}
\begin{center}
\begin{tikzpicture}
\begin{groupplot}[group style={group name=g, group size=3 by 1, horizontal sep=1.3cm}, ftdemof, height=0.18\textheight]
\nextgroupplot[title={(a) Espectro original\\[1pt]{\tiny\ftstrut $X_f(f)$}}, yticklabels={$\frac{1}{2}$}]
  \addplot[ftimpulso, ftazul!40!black] coordinates {(-\ftfzero,1) (\ftfzero,1)};
\nextgroupplot[title={(b) Trem de impulsos\\[1pt]{\tiny\ftstrut $P(f)$}}, yticklabels={$\frac{1}{T_s}$}]
  \addplot[ftimpulso, ftazul] coordinates {\ftimpf};
\nextgroupplot[title={(c) Sinal amostrado\\[1pt]{\tiny\ftstrut $X_\delta(f)$}}, yticklabels={$\frac{1}{2T_s}$}]
  \addplot[ftimpulso, ftazul] coordinates {\ftrepideal};
\end{groupplot}
\ftop{1}{2}{*}
\ftop{2}{3}{=}
\end{tikzpicture}
\end{center}
Todas as réplicas têm a mesma amplitude: não há envelope ponderando o espectro.
\end{frame}

\begin{frame}{Método Ideal: Amostragem}{Condição de Recuperação}
\small
A banda base ocupa $|f| \le f_{max}$ e a réplica vizinha começa em $f_s - f_{max}$.
Para que não haja sobreposição:
\begin{equation*}
  f_s - f_{max} > f_{max}
  \quad\Longrightarrow\quad
  \boxed{\; f_s > 2 f_{max} \;}
  \qquad
  \text{(banda de guarda } = f_s - 2f_{max})
\end{equation*}
\begin{itemize}
  \item Satisfeita a condição, as réplicas ficam separadas e a banda base pode ser
        isolada --- é o que torna a reconstrução possível.
  \item Violada a condição, as réplicas invadem a banda base: \textit{aliasing}, que
        nenhum filtro posterior desfaz.
  \item Por isso o filtro anti-aliasing atua \textbf{antes} da amostragem, limitando
        $f_{max}$ ao valor previsto no projeto.
\end{itemize}
\end{frame}

% ----------------------------------------------------------------------------
%  B. APLICAÇÃO AOS PARÂMETROS FIXOS
% ----------------------------------------------------------------------------
\begin{frame}{Método Ideal: Aplicação}{Parâmetros Fixos --- Domínio do Tempo}
    \small
    Com $x_f(t) = \cos(2\pi \cdot 2t)$, $f_s = 5$~Hz e $T_s = 0,2$~s:
    \begin{align*}
        x_\delta(t) &= x_f(t) \cdot \sum_{n=-\infty}^{\infty} \delta(t - nT_s)
        = \sum_{n=-\infty}^{\infty} x_f(nT_s)\, \delta(t - nT_s) \\
        &= \sum_{n=-\infty}^{\infty} \cos\!\left(2\pi \cdot 2 \cdot \frac{n}{5}\right) \delta\!\left(t - \frac{n}{5}\right)
        = \sum_{n=-\infty}^{\infty} \cos\!\left(\frac{4\pi n}{5}\right) \delta\!\left(t - \frac{n}{5}\right)
    \end{align*}

    \textbf{Amostras (o padrão se repete a cada 5 amostras, pois $f_0/f_s = 2/5$):}
    \[
        \begin{array}{c|ccccc}
            n & 0 & 1 & 2 & 3 & 4 \\ \hline
            x_f(nT_s) & 1 & -0{,}809 & 0{,}309 & 0{,}309 & -0{,}809
        \end{array}
    \]

    São $2{,}5$ amostras por período do sinal: o mínimo exigido por $f_s > 2f_{max}$.
\end{frame}

\begin{frame}{Método Ideal: Aplicação}{Parâmetros Fixos --- Domínio da Frequência}
    \small
    Substituindo $f_s = 5$~Hz na expressão do espectro amostrado:
    $$X_\delta(f) = f_s \cdot \sum_{k=-\infty}^{\infty} X_f(f - k f_s) = 5 \cdot \sum_{k=-\infty}^{\infty} X_f(f - 5k)$$

    \textbf{Análise das componentes:}
    $$
    \left.\begin{matrix}
    k=0: \pm 2\text{ Hz} \\
    k=1: 3 \text{ e } 7\text{ Hz} \\
    k=-1: -7 \text{ e } -3\text{ Hz} \\
    k=2: 8 \text{ e } 12\text{ Hz}
    \end{matrix}\right\} \text{componentes em } \pm 2, \pm 3, \pm 7, \pm 8, \pm 12, \dots
    $$

    Cada réplica tem peso $f_s \cdot \frac{1}{2} = 2{,}5$: \textbf{todas com a mesma
    amplitude}, pois não há pulso ponderando o espectro. A banda base ($2$~Hz) e a
    primeira réplica ($3$~Hz) estão separadas --- não há \textit{aliasing}.
\end{frame}

% ----------------------------------------------------------------------------
%  C. RECONSTRUÇÃO
% ----------------------------------------------------------------------------
\begin{frame}{Método Ideal: Reconstrução}{Filtro Passa-Baixas Ideal}
    \small
    Separadas as réplicas, basta descartar todas menos a banda base. O filtro de
    reconstrução é um passa-baixas ideal de corte $f_s/2$ e ganho $T_s$:
    \[
        H_{PB}(f) =
        \begin{cases}
            T_s, & |f| < \dfrac{f_s}{2}\\[4pt]
            0, & \text{caso contrário}
        \end{cases}
        \qquad\longrightarrow\qquad
        \text{com os parâmetros fixos: }
        \begin{cases}
            0{,}2, & |f| < 2{,}5\,\text{Hz}\\
            0, & \text{caso contrário}
        \end{cases}
    \]

    \textbf{Verificação da separação:}
    \[
        \underbrace{2\,\text{Hz}}_{\text{banda base}} < \underbrace{2{,}5\,\text{Hz}}_{\text{corte}} < \underbrace{3\,\text{Hz}}_{\text{1ª réplica}}
        \qquad \Longrightarrow \qquad
        \text{margem de } 0{,}5\,\text{Hz de cada lado}
    \]

    Essa margem é a banda de guarda de $1$~Hz obtida com $f_s = 5$~Hz: se
    $f_s \le 4$~Hz, a réplica em $f_s - 2$ entraria na banda base.
\end{frame}

\begin{frame}{Método Ideal: Reconstrução}{Demonstração no Domínio da Frequência}
\small
A reconstrução é o produto do espectro amostrado pelo passa-baixas. O ganho $T_s$
cancela exatamente o fator $f_s$ introduzido pela amostragem:
\begin{equation*}
  X_r(f) = X_\delta(f)\cdot \idrec{H_{PB}(f)} = T_s\cdot\frac{1}{T_s}X_f(f) = \boxed{\,X_f(f)\,}
  \qquad
  \text{(com os valores: } 0{,}2 \cdot 5 \cdot \tfrac{1}{2} = \tfrac{1}{2}\text{)}
\end{equation*}
\begin{center}
\begin{tikzpicture}
\begin{groupplot}[group style={group name=g, group size=3 by 1, horizontal sep=1.3cm}, ftdemof, height=0.18\textheight]
\nextgroupplot[title={(a) Espectro amostrado\\[1pt]{\tiny\ftstrut $X_\delta(f)$ e $H_{PB}(f)$}}, yticklabels={$\frac{1}{2T_s}$}]
  \draw[ftroxo, densely dashed, fill=ftroxo!8] (axis cs:-\ftfc,0) rectangle (axis cs:\ftfc,1.15);
  \addplot[ftimpulso, ftazul] coordinates {\ftrepideal};
\nextgroupplot[title={(b) Passa-baixas ideal\\[1pt]{\tiny\ftstrut $H_{PB}(f)/T_s$}}, yticklabels={$1$},
  visible on=<2->]
  \addplot[ftroxo, thick, fill=ftroxo!8, visible on=<2->] coordinates
    {(-\ftfc,0) (-\ftfc,1) (\ftfc,1) (\ftfc,0)};
\nextgroupplot[title={(c) Sinal reconstruído\\[1pt]{\tiny\ftstrut $X_r(f) = X_f(f)$}}, yticklabels={$\frac{1}{2}$}]
  \addplot[ftimpulso, ftverde, visible on=<3->] coordinates {(-\ftfzero,1) (\ftfzero,1)};
\end{groupplot}
\ftop{1}{2}{\times}
\ftop{2}{3}{=}
\end{tikzpicture}
\end{center}
Sem ponderação por $\mathrm{sinc}$ e sem termo de fase: o espectro recuperado é
\textbf{idêntico} ao original, dispensando equalizador.
\pause\pause
\end{frame}

\begin{frame}{Método Ideal: Reconstrução}{Demonstração no Domínio do Tempo}
\small
No tempo, a reconstrução é a convolução do sinal amostrado com a resposta ao impulso do
passa-baixas, $h_{PB}(t) = \mathrm{sinc}(f_s t)$. Cada impulso vira um $\mathrm{sinc}$
ponderado pela sua amostra:
\begin{equation*}
  x_r(t) = x_\delta(t) * \idrec{h_{PB}(t)}
  = \sum_{n=-\infty}^{\infty} x_f(nT_s)\, \idrec{\mathrm{sinc}\!\big(f_s(t - nT_s)\big)}
  = \boxed{\, x_f(t) \,}
\end{equation*}
A interpolação é exata porque $\mathrm{sinc}(f_s(t-nT_s))$ vale $1$ em $t = nT_s$ e
$0$ nos demais instantes de amostragem: as caudas se cancelam entre si.
\begin{center}
\begin{tikzpicture}
\begin{groupplot}[group style={group name=g, group size=3 by 1, horizontal sep=1.3cm}, ftdemot, height=0.18\textheight]
\nextgroupplot[title={(a) Sinal amostrado\\[1pt]{\tiny\ftstrut $x_\delta(t)$}}]
  \addplot[gray!55, dashed, domain=0:\fttmax, samples=300] {cos(deg(2*pi*\ftfzero*x))};
  \addplot[ycomb, ftazul, thick] coordinates {\ftamostras};
  \addplot[ftimppos, ftazul] coordinates {\ftamospos};
  \addplot[ftimpneg, ftazul] coordinates {\ftamosneg};
\nextgroupplot[title={(b) Núcleos de interpolação\\[1pt]{\tiny\ftstrut $x_f(nT_s)\,\mathrm{sinc}(f_s(t-nT_s))$}}, clip=true]
  \idnucleos
\nextgroupplot[title={(c) Soma dos núcleos\\[1pt]{\tiny\ftstrut $x_r(t) = x_f(t)$}}]
  \draw[black!45, densely dashed, thin] (axis cs:0,1) -- (axis cs:\fttmax,1);
  \addplot[gray!55, dashed, domain=0:\fttmax, samples=300] {cos(deg(2*pi*\ftfzero*x))};
  \addplot[ftverde, thick, densely dashed, domain=0:\fttmax, samples=300, visible on=<3->]
    {cos(deg(2*pi*\ftfzero*x))};
\end{groupplot}
\end{tikzpicture}\par\vspace{2pt}
{\scriptsize
\tikz[baseline=-0.6ex]\draw[ftazul, thick] (0,0) -- (0.4,0);\; amostras $x_f(nT_s)$
\qquad
\tikz[baseline=-0.6ex]\draw[ftroxo!55] (0,0) -- (0.4,0);\; núcleos $\mathrm{sinc}$
\qquad
\tikz[baseline=-0.6ex]\draw[ftverde, thick, densely dashed] (0,0) -- (0.4,0);\; sinal reconstruído
\qquad
\tikz[baseline=-0.6ex]\draw[gray!55, dashed] (0,0) -- (0.4,0);\; $x_f(t)$}
\end{center}
\pause\pause
\end{frame}

% ----------------------------------------------------------------------------
%  D. VISÃO GERAL DO PROCESSO + GANCHO
% ----------------------------------------------------------------------------
\begin{frame}{Método Ideal}{Visão Geral no Domínio do Tempo}
\centering
\begin{tikzpicture}
\begin{groupplot}[
  group style={group size=3 by 1, horizontal sep=1.2cm},
  ftaxis, width=0.26\textwidth
]

% (a) Sinal original
\nextgroupplot[title={(a) Sinal original\\[1pt]{\tiny\ftstrut $x_f(t)=\cos(2\pi\cdot 2t)$}}]
  \addplot[ftazul, thick, domain=0:\fttmax, samples=400]
    {cos(deg(2*pi*\ftfzero*x))};

% (b) Sinal amostrado (impulsos ponderados)
\nextgroupplot[title={(b) Sinal amostrado\\[1pt]{\tiny\ftstrut $x_\delta(t)$}}]
  \addplot[gray!55, dashed, domain=0:\fttmax, samples=400]
    {cos(deg(2*pi*\ftfzero*x))};
  \addplot[ycomb, ftazul, thick] coordinates {\ftamostras};
  \addplot[ftimppos, ftazul] coordinates {\ftamospos};
  \addplot[ftimpneg, ftazul] coordinates {\ftamosneg};

% (c) Após o passa-baixas: sinal idêntico ao original
\nextgroupplot[title={(c) Após o passa-baixas\\[1pt]{\tiny\ftstrut $x_r(t) = x_f(t)$}},
  clip mode=individual]
  \draw[black!45, densely dashed, thin] (axis cs:0,1) -- (axis cs:\fttmax,1);
  \addplot[gray!55, dashed, domain=0:\fttmax, samples=400]
    {cos(deg(2*pi*\ftfzero*x))};
  \addplot[ftverde, thick, domain=0:\fttmax, samples=400]
    {cos(deg(2*pi*\ftfzero*x))};
  \draw[gray, densely dotted] (axis cs:\ftpico,1) -- (axis cs:\ftpico,1.14);
  \node[font=\tiny, anchor=south west, inner sep=1pt] at (axis cs:\ftpico,1.12) {sem atraso};

\end{groupplot}
\end{tikzpicture}

\vspace{0.4em}
\mbox{\scriptsize
\tikz[baseline=-0.6ex]\draw[gray!55, dashed] (0,0) -- (0.45,0); $x_f(t)$\hspace{0.7em}
\tikz[baseline=-0.6ex]\draw[ftazul, thick] (0,0) -- (0.45,0); impulsos de área $x_f(nT_s)$\hspace{0.7em}
\tikz[baseline=-0.6ex]\draw[ftverde, thick] (0,0) -- (0.45,0); após o passa-baixas}

\vspace{0.4em}
{\small Amplitude preservada e nenhum atraso: não há equalizador a projetar.}
\end{frame}

\begin{frame}{Método Ideal}{Visão Geral no Domínio da Frequência}
\centering
\begin{tikzpicture}
\begin{groupplot}[
  group style={group size=3 by 1, horizontal sep=1.2cm},
  ftfreqaxis, width=0.26\textwidth
]

% (a) Espectro do sinal original
\nextgroupplot[title={(a) Espectro original\\[1pt]{\tiny\ftstrut $X_f(f)=\frac{1}{2}[\delta(f-2)+\delta(f+2)]$}}]
  \addplot[ftimpulso, ftazul] coordinates {(-\ftfzero,0.5) (\ftfzero,0.5)};

% (b) Espectro amostrado: réplicas em n f_s, todas com a mesma amplitude
\nextgroupplot[title={(b) Sinal amostrado\\[1pt]{\tiny\ftstrut $X_\delta(f)$}}]
  \draw[black!55, densely dashed] (axis cs:-\ftfmax,0.5) -- (axis cs:\ftfmax,0.5);
  \addplot[ftimpulso, ftlaranja!90!black] coordinates {\idrepmeio};
  \addplot[ftimpulso, ftazul] coordinates {(-\ftfzero,0.5) (\ftfzero,0.5)};

% (c) Após o passa-baixas: só a banda base, sem atenuação
\nextgroupplot[title={(c) Após o passa-baixas\\[1pt]{\tiny\ftstrut $X_r(f) = X_f(f)$}}]
  \draw[black!45, densely dashed, thin] (axis cs:-\ftfmax,0.5) -- (axis cs:\ftfmax,0.5);
  \draw[ftroxo, densely dashed, fill=ftroxo!8] (axis cs:-\ftfc,0) rectangle (axis cs:\ftfc,0.6);
  \node[font=\tiny, anchor=south, inner sep=1pt, yshift=4pt, ftroxo] at (axis cs:0,0.6) {$H_{PB}(f)$};
  \addplot[ftimpulso, ftverde] coordinates {(-\ftfzero,0.5) (\ftfzero,0.5)};

\end{groupplot}
\end{tikzpicture}

\vspace{0.4em}
\mbox{\scriptsize
\tikz[baseline=-0.6ex]\draw[black!55, densely dashed] (0,0) -- (0.45,0); envelope constante $\frac{1}{2}$\hspace{0.7em}
\tikz[baseline=-0.6ex]\draw[ftazul, thick] (0,0) -- (0.45,0); banda base\hspace{0.7em}
\tikz[baseline=-0.6ex]\draw[ftlaranja!90!black, thick] (0,0) -- (0.45,0); réplicas em $nf_s \pm 2$\hspace{0.7em}
\tikz[baseline=-0.6ex]\draw[ftroxo, densely dashed, fill=ftroxo!8] (0,-0.08) rectangle (0.5,0.12);\hspace{3pt}passa-baixas}

\vspace{0.4em}
{\small Envelope constante: todas as réplicas com a mesma amplitude, portanto nenhuma
distorção na banda base.}
\end{frame}
